% =======================================
% MEETING 8
% ---------------------------------------
% DATE   - September 16, 2026
% FORMAT - F2F 
% =======================================
% TOPIC  - Dual Vectors
% =======================================

\subsection{Dual Vectors}

	From the last lecture, we learned that a vector field is a selection
	of a vector for each tangent space $T_p \mcal M$, i.e., we select one
	vector $v$ for each vector space, which is present for each point in the 
	manifold. We call the union of these tangent spaces for all of $p$ to be
	a tangent bundle $T\mcal M$.
	\[
		T\mcal M = \bigcup_p T_p \mcal M
	\]
	and a vector field $v$ is a section of that tangent bundle.
	\[
		v \in T\mcal M
	\]
	Now, recall what the action of a vector at a point is; It is a derivation that can
	act on functions, which returns back a real number $v_p: \mcal F \to \R$. i.e., $v_p[f]\in\R$. It 
	is a map which gets fed a function, or rather a set of functions for each 
	coordinate, and it returns the rate of change for that particular direction, which is 
	\[
		v_p:\mcal{F} \to \R
	\]

	Similarly, a vector field is another operation or derivation, which gets fed a set
	of functions, and returns another set of functions. Instead of getting a set of numbers, we 
	receive another set of functions.
	\[
		v:\mcal{F}\to\mcal{F}
	\]

	Recall a scalar field. A scalar field is a simply a smooth function from the manifold to the real numbers.
	\[
		f:\mcal M \to \R,\qquad f\in\mcal{F}
	\]
	At each point $p\in\mcal M$, a scalar field assigns a number:
	\[
		f(p)\in\R
	\]

	A tangent vector $v_p\in T_p\mcal M$ acts on a scalar field by taking its directional
	derivative at $p$, i.e.,
	\[
		v_p[f]\in\R
	\]

	A vector field $v$ then assigns a tangent vector $v_p$ to every point $p$ on a manifold, so its action
	on a scalar field produces another scalar field, $v:\mcal{F}\to\mcal{F}$, that is:
	\[
		(v[f])(p) = v_p[f]
	\]

	As an example consider the real line $\R$. This is a manifold. For each point $p\in\R$, there 
	exists a tangent space $T_p\R$. Vectors exist on this tangent space. However, it is usually 
	conventional to represent vectors themselves as existing on the real line (refer to Fig. \ref{fig:vectorfieldRtrad}), while on reality,
	they exist on the tangent space. This is coincidental, because $\R$ and $T_p\R$ are isomorphic, i.e., $\R\simeq T_p\R$. However,
	if we are to be general, this is confusing, and we lose sight of the distinction between the tangent space
	and the manifold. We can instead represent $v$ as a function on the real line, and represent a vector as the values of these functions.
	\begin{figure}[h]
		\centering
		\includegraphics[width=4in, height=2in]{images/vectorsasfunctions.pdf}
		\caption{defaultcaption} % to be filled
		\label{fig:defaultlabel}
	\end{figure}

	In this illustration, then, we can imagine that $\cup_p T_p\R$ is a plane. 
	Most of this, we discussed last meeting. 
	Recall that a tangent vector acts only on a point. It get's fed a function, and evaluates
	the \textit{rate of change} of that function at that point only, and returns a number. 
	A vector field \textit{acts} on all the points defined on the manifold, and assigns a vector there. 

	We can write it like so: For a chart $\psi$ with coordinate functions
	$\{x^\nu\}$, we can represent a vector field as
	\begin{equation}
		\label{eq:vecfield}
		v = v^\mu(x^\nu)\pdv{}{x^\mu}.
	\end{equation}
	We write $v^\mu(x^\nu)$ because the components of the vector field depend
	on the choice of coordinates. Each $v^\mu$ is itself a scalar field.

	Recall that a coordinate function takes a point $p$ and assigns it a
	real number. Thus, $x^\sigma$ takes the point $p$ and returns its
	$\sigma$-th coordinate:
	\[
		x^\sigma:p\mapsto x^\sigma(p).
	\]

	A vector field can act on these coordinate functions. In particular,
	\begin{align*}
		v[x^\sigma]
		&= v^\mu(x^\nu)\pdv{x^\sigma}{x^\mu} \\
		&= v^\mu(x^\nu)\tensor{\delta}{^\sigma_\mu} \\
		&= v^\sigma(x^\nu).
	\end{align*}
	Therefore, when a vector field acts on the $\sigma$-th coordinate
	function, it returns the $\sigma$-th component of the vector field:
	\[
		v[x^\sigma]=v^\sigma.
	\]

	For example, in three dimensions,
	\[
		v =
		v^1(x^1,x^2,x^3)\pdv{}{x^1}
		+
		v^2(x^1,x^2,x^3)\pdv{}{x^2}
		+
		v^3(x^1,x^2,x^3)\pdv{}{x^3}.
	\]
	Each $v^\mu$ is a scalar field in its own right. Acting on the
	coordinate functions extracts these scalar fields:
	\[
		v[x^1]=v^1,\qquad
		v[x^2]=v^2,\qquad
		v[x^3]=v^3.
	\]
	Each $v^\mu$ is a scalar field in of itself. 

	\subsubsection{Dual Space}

		Here, we further discuss what we introduced in the previous chapters. The dual space, 
		also known as the covector space, or covariant space, or cotangent space. It has a lot of names, we will refer to covectors, duals, covariant vectors interchangably.

		Recall that a vector is a derivation, which acts on functions and returns real numbers.
		Duals, then, are algebraic structures which act on a vector and returns a real number.
		From this, we can define the dual space as the following:

		\begin{definition}{Dual Space}
			For a vector space $V$ with an element $v\in V$, a dual space
			$V^*$ is a space of linear functionals on $V$, such that a \textit{dual} $\omega\in V^*$
			acts on $v$ with the following property:
			\[
				\omega:V \to \R,\qquad v\mapsto \omega(v)
			\]
		\end{definition}
		We can then distinguish a tangent vector, whose input is a scalar field or function
		\[
			v_p:\mcal F \to \R
		\]
		whereas, for a  dual vector, the input is a tangent vector.
		\[
			\omega_p : \TpM \to \R
		\]

		In Physics, particularly in Quantum Theory, we have the \textit{bra} operator, $\bra{\cdot}$. This 
		is a covector. In particular, for functions $\alpha$ and $\beta$, we have
		\[
			\braket{\alpha}{\beta}\in\mathbb{C}
		\]
		This already has the property of linearity we know
		\[
			\bra{\alpha}\pqty{
				A\ket{\xi} + B\ket{\sigma}
			} = 
			A\braket{\alpha}{\xi} + B \braket{\alpha}{\sigma}
		\]
		Before QM, this is also already encountered in linear Algebra. If you consider
		a column matrix to be a vector, then a row matrix is a covector. They are homeomorphic
		to vectors and covectors in differential geometry.

		If $\displaystyle \bmqty{a\\b\\c}\in V$, then $\displaystyle \bmqty{
		d & e & f}\in V^*$. Notice, then, that when you take their product, the result
		is a real number:
		\[
			\bmqty{
				d&e&f
			} 
			\bmqty{
				a\\b\\c
			} 
			=
			ad + be + cf \in \R
		\]
		The only restriction we find ourselves in GR is that we are limited
		to the study of real numbers for this course, for majority of our encounters
		with physical backing all take place in the real axis. There exists a branch of mathematics 
		called complex differential geometry, it is an interesting field, but for most of our purposes, we shall
		not need it.

		The set of all dual vectors happen to be a vector space as well. To define 
		a dual vector you need to provide specific instructions on how it acts
		on all vectors. Once you have done that, you have defined a covector completely.

		\paragraph{Notation.} From here on out, we shall sporadically denote a vector to have an overline, $\overline{v}$, and 
		a dual to have an underline, $\underline{\omega}$. This is purely for teaching purposes, as it allows us to 
		identify which is which immediately. However, after some time, when you are already 
		able to identify duals and vectors within proper context, we will eventually rid ourselves of this pedagogical tool.

	\subsubsection{Properties and Operations}

		For the dual space to be a vector space, it must satisfy the 
		conditions we set to define a vector space, and that is linearity,
		multiplicative scaling, and the existence of a zero vector.

		As an operator, we require the following properties for a dual vector.
		\begin{enumerate}
			\item Linearity

				For $\underline{\omega}_1,\underline{\omega}_2\in V^*$ and $\overline{x}\in V$,
				\[
					\pqty{
						\underline{\omega}_1
						+
						\underline{\omega}_2
					}
					[\overline{x}]
					:=
					\underline{\omega}_1[\overline{x}]
					+
					\underline{\omega}_2[\overline{x}]
				\]
				
			\item Scalar multiplication

				For all $\alpha\in\R$, $\overline x \in V$ and $\underline{\omega}\in V^*$,
				\[
					\pqty{
						\alpha
						\underline{\omega}
					}
					[\overline{x}]
					:=
					\alpha\underline{\omega}[\overline{x}]
				\]

			\item Zero covector

				There exists a unique covector $\underline{0}\in V^*$ such that, for all $\overline{x}\in V$,
				\[
					\underline{0}[\overline{x}] = 0
				\]
		\end{enumerate}

		An operator on $v\in V$ must satisfy these conditions before they 
		can be considered to belong in $V^*$.

	\subsubsection{Dual Basis}

		Recall that the dual space $V^*$ of a vector space $V$ is the space of all
		linear maps from $V$ to $\R$:
		\[
			V^*
			=
			\left\{
				\underline{\omega}:V\to\R
				\;\middle|\;
				\underline{\omega}\text{ is linear}
			\right\}.
		\]
		Thus, a covector does not require a manifold to be defined. On a manifold,
		the corresponding construction is applied to the tangent space at a point.
		A covector is instead a map that defines an operation to vectors. It is an 
		abstract dual vector space. However, a dual vector space needs a basis, and as such,
		we define the following:

		\begin{definition}{Dual basis}
			Let $\{\overline{e}_\mu\}$ be a basis for $V$. We define the corresponding
			dual basis $\{\underline{e}^\mu\}$ of $V^*$ that acts on the basis of $V$
			such that
			\begin{equation}
				\label{eq:dual-basis-definition}
				\underline{e}^\mu[\overline{e}_\nu]
				=
				\tensor{\delta}{^\mu_\nu}.
			\end{equation}
		\end{definition}

		That is, the $\mu$-th covector returns $1$ when acting on the $\mu$-th
		basis vector and $0$ when acting on every other basis vector.

		Since $\{\overline{e}_\mu\}$ is a basis for $V$, an arbitrary vector
		$\overline{v}\in V$ can be written as
		\[
			\overline{v}
			=
			v^\nu\overline{e}_\nu.
		\]
		\begin{proof}
			Using the linearity of $\underline{e}^\mu$, its action on $\overline{v}$
			is therefore
			\begin{align}
				\underline{e}^\mu[\overline{v}]
				&=
				\underline{e}^\mu
				\left[
					v^\nu\overline{e}_\nu
				\right]
				\nonumber\\
				&=
				v^\nu
				\underline{e}^\mu[\overline{e}_\nu]
				\nonumber\\
				&=
				v^\nu
				\tensor{\delta}{^\mu_\nu}
				\nonumber\\
				&=
				v^\mu.
				\label{eq:dual-action}
			\end{align}
			Hence, the action of the $\mu$-th dual basis vector extracts the $\mu$-th
			component of $\overline{v}$.
		\end{proof}

		We now show that $\{\underline{e}^\mu\}$ is a basis for $V^*$ by proving
		linear independence and spanning.

		\paragraph{Linear independence}

		\begin{proof}
			Suppose that a linear combination of the dual basis vectors vanishes:
			\[
				a_\mu\underline{e}^\mu
				=
				\underline{0},
			\]
			where $a_\mu\in\R$. Acting on an arbitrary basis vector
			$\overline{e}_\nu$ gives
			\begin{align}
				0
				&=
				\underline{0}[\overline{e}_\nu]
				\nonumber\\
				&=
				\left(
					a_\mu\underline{e}^\mu
				\right)
				[\overline{e}_\nu]
				\nonumber\\
				&=
				a_\mu
				\underline{e}^\mu[\overline{e}_\nu]
				\nonumber\\
				&=
				a_\mu
				\tensor{\delta}{^\mu_\nu}
				\nonumber\\
				&=
				a_\nu.
			\end{align}
			Since this holds for every $\nu$, we have
			\[
				a_\nu=0
			\]
			for every $\nu$. Therefore,
			\[
				\left\{
					\underline{e}^\mu
				\right\}
			\]
			is linearly independent.
		\end{proof}

		\paragraph{Spanning}

		\begin{proof}
			Let $\underline{\omega}\in V^*$ be an arbitrary covector. We want to
			show that $\underline{\omega}$ can be written as a linear combination of
			the dual basis vectors.

			Suppose we write the components of a covector as $\omega_\mu$. Consider the arbitrary 
			covector
			\[
				\underline{\omega}'
				=
				\omega_\mu\underline{e}^\mu.
			\]
			Acting on an arbitrary vector $\overline{v}=v^\nu\overline{e}_\nu$ gives
			\begin{align}
				\underline{\omega}'[\overline{v}]
				&=
				\left(
					\omega_\mu\underline{e}^\mu
				\right)
				[\overline{v}]
				\nonumber\\
				&=
				\omega_\mu
				\underline{e}^\mu[\overline{v}]
				\nonumber\\
				&=
				\omega_\mu v^\mu.
			\end{align}
			On the other hand, by the linearity of $\underline{\omega}$,\footnote{We shall prove this relation in a bit}
			\begin{align}
				\underline{\omega}[\overline{v}]
				&=
				\underline{\omega}
				\left[
					v^\mu\overline{e}_\mu
				\right]
				\nonumber\\
				&=
				v^\mu
				\underline{\omega}[\overline{e}_\mu]
				\nonumber\\
				&=
				v^\mu\omega_\mu.
			\end{align}
			Therefore,
			\[
				\underline{\omega}'[\overline{v}]
				=
				\underline{\omega}[\overline{v}]
			\]
			for every $\overline{v}\in V$. Hence,
			\[
				\underline{\omega}'
				=
				\underline{\omega},
			\]
			and consequently
			\begin{equation}
				\label{eq:covector-expansion}
				\underline{\omega}
				=
				\omega_\mu\underline{e}^\mu,
				\qquad
				\omega_\mu
				=
				\underline{\omega}[\overline{e}_\mu].
			\end{equation}
			Thus, the dual basis spans $V^*$.
		\end{proof}

		Since $\{\underline{e}^\mu\}$ is both linearly independent and spanning,
		it is indeed a basis for $V^*$. 

		Note the following: We can extract the components of a covector by feeding it 
		a basis vector. Consider
		\begin{align*}
			\underline{\omega}
			[\overline{e}_\nu]
			&=
			\omega_\mu \underline{e}^\mu
			[\overline{e}_\nu] \\
			&= \omega_\mu \tensor{\delta}{^\mu_\nu}\\
			&= \omega_\nu
		\end{align*}
		That is, feeding the $\nu$-th basis vector on a dual vector will get us the 
		$\nu$-th component of the covector.

		We can therefore write vectors and covectors as the following
		\begin{align*}
			\underline{\omega} = \omega_\mu \underline{e}^\mu &&
			\overline{v} = v^\mu \overline{e}_\mu
		\end{align*}
		These are fundamentally different abstract geometric objects. They just happen to be isomorphic.
		One acts on functions, the other acts on vectors. 

	\subsubsection{Double-dual space}
		Consider: Since $V^*$ is a vector space, and a dual vector space acts on a vector space, shouldn't there be a dual space for a dual vector space? 
		There should be, and in fact this chain can be continued ad infinitum:
		\[
			V \to V^* \to V^{**} \to \cdots \to V^{**\cdots**}.
		\]

		However, there is a natural relationship between $V$ and its double dual $V^{**}$. For any $v\in V$, we can 
		define an element of $V^{**}$ by
		\[
			\uuline{\Omega} [\underline \omega] := \underline\omega[\overline v],
			\qquad \underline\omega\in V^*.
		\]
		Thus, there is a natural map
		\[
			V \to V^{**},
			\qquad
			\overline v \mapsto \uuline{\Omega}
		\]
		This map is always injective. If $V$ is finite-dimensional, 
		it is also surjective, and hence
		\[
			V^{**}\cong V.
		\]
		Thus, for finite-dimensional vector spaces, taking the 
		dual twice brings us back to a space naturally isomorphic 
		to the original vector space.

		That is, if $\uuline{\Omega}\in V^{**}$, then
		\[
			\uuline{\Omega}[\underline\omega]\in\mathbb R,
			\qquad \underline\omega \in V^*.
		\]
		When $V$ is finite-dimensional, every such $\uuline\Omega$ can be 
		identified with some $v\in V$, such that
		\begin{equation}
			\label{eq:double-dual-vector-congruency}
			\uuline{\Omega}[\underline\omega]=\underline\omega[\overline v].
		\end{equation}
		In this sense, $V$ and $V^{**}$ are naturally identified and correspond with one another.

		\paragraph{Going back to manifolds}

		On a manifold, vector spaces are replaced by tangent spaces
		$T_p\mathcal M$. Their dual spaces are the cotangent spaces
		\[
			T_p^*\mathcal M := (T_p\mathcal M)^*.
		\]
		This is where covectors, or cotangent vectors, live.

	\subsubsection{The Differential of a Function}

		Now that we have characterized duals abstractly, we can define a particularly important covector: the differential of a function.

		\begin{definition}{Differential of a function}
			Given $f\in\mathcal F(\mathcal M)$, its differential at $p\in\mathcal M$ is the covector
			\[
				(\underline{df})_p\in T_p^*\mathcal M
			\]
			defined by
			\[
				(\underline{df})_p:T_p\mathcal M\to\mathbb R,
				\qquad
				\overline{v}_p\mapsto (\underline{df})_p[\overline{v}_p]:=\overline{v}_p[f].
			\]
		\end{definition}

		That is, the differential of a function takes a tangent vector and returns the directional derivative of the function in that direction. In other words,
		\[
			(\underline{df})_p[\overline{v}_p]=\overline{v}_p[f].
		\]
		Since $\overline{v}_p[f]$ is linear in $\overline v_p$, $(\underline{df})_p$ is a linear map and therefore a covector.

		If we define the differential at every point $p\in\mathcal M$, we obtain a covector field:
		\[
			\underline{df}\in T^*\mathcal M.
		\]
		Thus,
		\[
			\underline{df}\iff\text{a covector field}.
		\]

		Covector fields are called \textit{one-forms}. They are sections of the cotangent bundle $T^*\mathcal M$, where
		\[
			T^*\mathcal M
			=
			\bigcup_{p\in\mathcal M}T_p^*\mathcal M.
		\]

		To put it more physically, Lagrangian mechanics naturally lives on the tangent bundle $T\mathcal M$, whose points specify positions and velocities. Hamiltonian mechanics naturally lives on the cotangent bundle $T^*\mathcal M$, whose points specify positions and momenta. The latter is the phase space of Hamiltonian mechanics.

		\subsubsection{$n$-forms}

		An $n$-form is defined as a wedge product of $1$-forms. 

		A wedge product (or \textit{exterior product}) is an operation which
		combines vectors, or more commonly, covectors/forms, to produce an antisymmetric 
		higher-rank objects.

		For two 1-forms $\alpha, \beta \in V^*$, their wedge-product is a $2$-form:
		\[
			\alpha \wedge \beta
		\]
		It is defined by antisymmetrizing the tensor product (we will discuss what it is in \S \ref{subsec:tensor product}). That is,
		\[
			\alpha\wedge\beta = \alpha\otimes\beta - \beta\otimes\alpha
		\]
		so that, acting on two vectors $u, v$,
		\[
			(\alpha \wedge \beta)[u,v] = \alpha[u]\beta[v]-\alpha[v]\beta[u]
		\]
		It is antisymmetric because if we exchange the two arguments,
		\[
			(\alpha \wedge \beta)[v,u] = \alpha[v]\beta[u]-\alpha[u]\beta[v]
		\]
		so
		\[
			(\alpha \wedge \beta)[v,u] = - (\alpha\wedge\beta)[u,v]
		\]

		In particular, 
		\begin{equation}
			\label{eq:self-wedge}
			\alpha\wedge\alpha = 0
		\end{equation}

		An $n$-form is a basis for geometric integration, and wedge products gives us a natural
		langauge for integration on manifolds. Say for example the coordinate 1-forms, $\dd{x}^1$ and $\dd{x}^2$.
		Let $x^1 =x$ and $x^2 = y$. When we integrate an area element, we have
		\[
			\int f \dd{A} = \iint f \dd{x}\dd{y}
		\]
		Properly speaking, $\dd{x}\dd{y}$ should be a two-form, $\dd{x}\wedge\dd{y}$. It is an oriented area element. Therefore, what we really
		ought to write is
		\[
			\int f \dd{x}\wedge \dd{y}
		\]

		Similarly, $\dd{V} = \dd{x}\dd{y}\dd{z}$ is properly $\dd{x}\wedge\dd{y}\wedge\dd{z}$, a three-form. It is an oriented volume element. 

		Therefore, when we see expressions such as 
		\[
			\int\omega
		\]
		This means we are integrating $n$-forms. A differential, and the wedge product of 
		differentials, are the basis for integration.

		More generally, if $\omega$ is a $k$-form and $\eta$ is an $\ell$-form,
		then their wedge product is a $(k+\ell)$-form:
		\[
			\omega\wedge\eta\in\Lambda^{k+\ell}(V^*).
		\]
		Thus, we can construct $n$-forms by taking wedge products of $1$-forms,
		such as
		\[
			\dd{x}^1\wedge\dd{x}^2\wedge\cdots\wedge\dd{x}^n
		\]

		On an $n$-dimensional manifold, an $n$-form can be integrated over an
		$n$-dimensional oriented region:
		\[
			\int_{\mathcal M}\omega
		\]
		More generally, a $k$-form is integrated over a $k$-dimensional oriented
		manifold or submanifold.

		Therefore, differential forms provide the mathematical objects that
		underlie the familiar notation of integration. The symbols $\dd{x}$,
		$\dd{x}\wedge\dd{y}$, and
		$\dd{x}\wedge\dd{y}\wedge\dd{z}$ should not merely be thought of as
		infinitesimal quantities; they are, respectively, $1$-, $2$-, and
		$3$-forms.

		This is not very important for the beginning of GR, but it will become
		very useful in the latter parts, when deriving GR using the variational principle.

		In elemetary calculus and real anaylsis, we naturally think of this \textit{differential} as a small change, an \textit{infintesimal}. 
		\[
			x + \dd{x} = \text{a small change in }x
		\]
		If mathematicians saw you using that literally, they might throw a book at you. This is not at all what 
		we mean by a differential, this is a heuristic. A very useful heuristic, sure, but it's wrong.
		It's only a shorthand for thinking about a quantity that is arbitrarily small. What we mean when we say $dx$ is that it's a limit as $\Delta x$ vanishes,
		but it's not equal to $0$ itself.
		\[
			x + dx = \lim_{\Delta x \to 0} x + \Delta x \quad \neq x
		\]
		We don't actually need this to perform integration, we can integrate with $n$ forms entirely, but it's a heuristic, and
		it's a pedagological tool, which allows us to think of things in an intuitive manner.
		Learning elementary calculus using the complexities of one-forms is a
		difficult endeavor. It is doable, yes, but you need to have learned differential calculus prior,
		which is pointless when you're trying to learn introductory calculus in the first place.

		Heuristics such as these make teaching and understanding things simpler.
		Heuristics are tools, and they are very useful tools. Do not throw your tools away for
		the sake of rigor. However, you must remember how these tools are justified, and what they
		really are under the hood, mathematically, lest you may lose sight of the operations you do when you extend 
		your calculations beyond which those heuristics apply.

	\subsubsection{Differential Basis}

		Moving forward, consider a chart with basis functions $\{x^\mu\}$. Recall that each $x^\mu$ is a
		scalar function. The differential of each of these functions is then
		\[
			dx^\mu
		\]
		Let's consider the following set of $n$ one-forms of the coordinate basis functions, 
		$\{\underline{\dd{x}}^\mu\}$.
		Let this act on a basis vector. We know that if you have a chart, the most natural
		basis vectors is the coordinate basis vector, hence:
		\begin{align*}
			\dd{x^\mu}\qty[
				\pdv{}{x^\nu}  
				] 
			&= \pdv{}{x^\nu}[x^\mu]\\
			&= \pdv{x^\mu}{x^\nu}\\
			&= \tensor{\delta}{^\mu_\nu}
		\end{align*}

		This is precisely the property of a dual basis, (\ref{eq:dual-basis-definition}). Hence
		$\dd{x^\mu}$ is a basis for $T^* \mcal M$. 

		Since $\underline{dx}^\mu$ is a basis, we get to write an arbitrary one-form as a linear 
		combination of components $\omega_\mu$ and the basis:
		\[
			\underline\omega = \omega_\mu \dd{x^\mu}
		\]
		and we also find that $\overline{v}[x^\mu]= \underline{dx}^\mu [v]$, as per the previously 
		defined properties. Since $\overline{v}[x^\mu] = v^\mu$, then $\underline{dx}^\mu [v] = v^\mu$ as previously shown. We can prove
		it explicitly, however.
		\begin{proof}
			\begin{align*}
				dx^\mu[v] 
				&= v[x^\mu]\\
				&= v^\nu \pdv{}{x^\nu} [x^\mu] \\
				&= v^\nu \pdv{x^\mu}{x^\nu}\\
				&= v^\nu \tensor{\delta}{^\mu_\nu}\\
				&= v^\mu
			\end{align*}
		\end{proof}

		Consider, that since we can write a one-form as 
		\[
			\underline{df} = \omega_\mu\dd{x^\mu}
		\]
		We must find an expression of the components $\omega_\mu$. We can do that
		by the following:

		Consider the differential at point $p$. 
		We can get the differential dual of a function $df$ for each point 
		on the manifold:
		\[
			(\underline{df})_p [v] = \overline{v}_p[f]
		\]
		We can then do a point by point substitution, and we have the covector field.\footnote{This 
		is also what we do to find the vector field, define it for all points in $\mcal M$}
		\[
			\underline{df}[v] = \overline{v}[f]
		\]
		We can expand $\overline v$ and calculate it over
		\begin{align*}
			\underline{df}[v] &= \overline{v}[f]\\
							  &= v^\mu \pdv{}{x^\mu}[f]\\
							  &= v^\mu \pdv{f}{x^\mu} \\
							  &= \dd{x^\mu}[v] \pdv{f}{x^\mu}
		\end{align*}
		However, $v$ is an arbitrary vector. In an operational sense, it is not needed. Hence, we have
		\[
			\underline{df} = dx^\mu \pdv{f}{x^\mu} 
		\]
		Since these are all just numbers, we can rearrange this, and we find the result is:
		\begin{equation}
			\label{eq:df-basis}
			\underline{\dd{f}} = \pdv{f}{x^\mu} \dd{x^\mu} 
		\end{equation}
		That is the components of a differential of a function $f$. However, recall that we 
		have seen this before. In classical real analysis, when we are asked to find the differential
		of $f(x,y,z)$, we write:
		\[
			df = \pdv{f}{x}dx + \pdv{f}{y}\dd{y} + \pdv{f}{z}\dd{z}  
		\]
		This is the same thing! Just reconceptualized. Since $\dd{x^\mu}$ are one forms, then,
		an arbitrary one form is simply a linear combination of other one-forms. 
		Given a chart, the differential of 
		a function is a one form whose components are the gradient and the basis are the differential
		of the coordinate basis functions. 

		That is, $\omega_\mu = \partial_\mu f$ and $e^\mu = dx^\mu$. Therefore
		\[
			df = \p_\mu f \dd{x^\mu}
		\]
		From this property, we can also define a change of basis, from $x^\mu$ to $x^\nu$. This is a dual transformation rule.
		\begin{equation}
			\label{eq:dual-change-basis}
		    dx^\nu = \pdv{x^\nu}{x^\mu} dx^\mu 
		\end{equation}
		
		We need a useful expression for a component of a covector given a basis. 
		We express an arbitrary covector as 
		\[
			\underline\omega = \omega_\mu dx^\mu
		\]
		Let's have it act on a coordinate basis vector $\p_\nu$. We have
		\begin{align*}
			\underline\omega[\p_\nu] 
			&= \omega_\mu \dd{x}^\mu[\p_\nu]\\
			&= \omega_\mu \tensor{\delta}{^\mu_\nu}\\
			&= \omega_\nu
		\end{align*}

		That is, if you want the $\nu$-th component of a one-form, you need to 
		feed it the $\nu$-th basis vector

	\subsubsection{The Interior Product}
		We can write a vector as
		\[
			\overline v = v^\mu \partial_\mu.
		\]
		We can write the component of the vector as an extraction of the dual basis
		\[
			v^\mu = \underline{dx}^\mu[v].
		\]
		However, from (\ref{eq:double-dual-vector-congruency}), recall that the double dual $V^{**}$ is canonically isomorphic to
		$V$, i.e., $V^{**}\cong V$, so that
		\[
			\uuline{v}[\underline{\omega}] = \underline{\omega}[\overline{v}].
		\]
		Therefore, since the two are isomorphic, we \emph{define} the following
		operation as equivalent:
		\[
			v^\mu = \overline{v}[\underline{dx}^\mu].
		\]
		This is called the \textit{interior product} (or \emph{contraction}).
		That is, we define
		\begin{align*}
			dx^\mu[v] &= v[dx^\mu] \\
					  &= i_\nu\, dx^\nu.
		\end{align*}
		This is an identity between the two ways of evaluating the pairing.

		We see that we have this universal feature: if we need the $\mu$-th
		component of an object, we feed it the $\mu$-th basis object. The
		expression $v[dx^\mu]$ is understood to mean ``evaluate the
		double-dual image of $v$ on the covector $dx^\mu$'' — and the result
		is the same number $v^\mu$ as $dx^\mu[v]$.

		\begin{theorem}{Component extraction}
			For a covector $u = u_i e^i$ and the dual basis $\{e^i\}$
			satisfying $e^i[e_j] = \tensor{\delta}{^i_j}$, the components are given
			by the standard pairing
			\[
				u_i = u[e_i].
			\]
			For a vector $v = v^i e_i$, the same components are recovered
			by evaluating the double-dual image of $v$ on the covectors
			$e^i$:
			\[
				v^i \equiv v[e^i] = e^i[v],
			\]
			where $v$ on the right is understood as an element of $V^{**}$
			via the canonical embedding.
		\end{theorem}

		Bases need not come from a chart, it is not necessary. A chart produces a
		\emph{coordinate basis}, but any set of linearly independent vectors
		spanning the tangent space is a legitimate basis. Given a chart, we
		can always write a general basis in terms of the coordinate basis,
		\[
			e_i = e_i{}^\mu \pdv{}{x^\mu},
		\]
		where $e_i{}^\mu$ are the components of the $i$-th basis vector.
		As long as the matrix $\pqty{e_i{}^\mu}$ has nonzero determinant,
		the vectors $\{e_i\}$ are linearly independent and hence form a basis:
		\[
			\det\pqty{e_i{}^\mu} \neq 0.
		\]
		A nonzero determinant is a requirement. 
		If the determinant were zero, the vectors would be
		linearly dependent and could not span the tangent space, so the whole framework
		collapses.

		There's a lot more to differential geometry 
		than what we currently show, a lot more fields than
		the Riemmannian that we study. Differential geometry is something that has been studied for over
		300 years, long before relativity. However, we can only teach you 
		the elements of differential geometry that is necessary to learn General Relativity.

		For theorists, it is natural to formulate physics in the language of
		differential geometry. In practice, however, computations are almost
		always carried out in coordinates: for example, the Schwarzschild
		solution is most conveniently written in the chart $(t, r, \theta, \phi)$.
		The basis vectors used in such computations are not necessarily tied to
		those coordinates — a point we will return to — but a coordinate basis
		is usually the simplest to work with, and from it one can, in principle,
		free oneself from coordinates entirely.

		Observers, i.e., experimental physicists, however, do not have access to this abstract description.
		An experimentalist measures \emph{components}, and only components. A
		clock measures a proper time interval; a rod measures a proper length;
		a detector measures an energy or a momentum. These are all scalars
		obtained by contracting tensorial objects with the observer's own
		local frame — never the abstract tensors themselves.
		This raises a practical problem. Tensors in a general curved spacetime
		are defined at each point, but they live in the tangent space, and there
		is no canonical way to compare or evaluate them without introducing
		extra structure. What the observer actually carries with them is a
		\emph{local inertial frame}: a set of four orthonormal basis vectors
		attached to their worldline, one timelike (their 4-velocity) and three
		spacelike (their local spatial axes). This frame is what a physical
		observer uses to decompose any tensor into measurable components.

		The mathematical object that captures this is the \textit{orthonormal tetrad}
		(or \emph{vierbein}): a set of four orthonormal basis vectors defined at each
		point, consisting of one timelike vector corresponding to the observer's 4-velocity
		and three spacelike vectors corresponding to their local spatial axes. By projecting
		tensors onto this local frame, one obtains the components that the observer measures.
		The tetrad therefore bridges the abstract geometric description of spacetime and
		concrete measurements: it is defined geometrically independently of any particular
		coordinate system, while its components in a coordinate basis provide the quantities
		used in practical calculations.

		For us to understand what that all means, we shall begin to discuss tensors.
\newpage
\subsection{Tensors}
	Tensors are algebraic objects, similar to what we already defined before 
	for vectors and duals. However, it is more general than that.
	\subsubsection{Definition}
		\begin{definition}{Tensor}
			A tensor of $(k,l)$ type is a multilinear map 
			from the product of dual and vector spaces to real numbers,
			such that
			\begin{equation}
				\label{eq:tensor-defn}
				T:
				V^*_1 \times 
				V^*_2 \times 
				\cdots
				V^*_k \times
				V_1 \times
				V_2 \times
				\cdots
				V_l
				\to \R
			\end{equation}
		\end{definition}

		That is, you need to supply a tensor $k$ covectors and $l$ vectors 
		for it to map to a real number. A tensor is therefore a function, that accepts 
		covectors and vectors,
		with several \textit{slots} for each. It has $k$ slots for covectors, and $l$ slots
		for vectors.
		\[
			T(
			\underbrace{
				\underline{\phantom{x}},
				\underline{\phantom{x}},
				\underline{\phantom{x}},
				\ldots,
				\underline{\phantom{x}},
				\underline{\phantom{x}},
				\underline{\phantom{x}}
			}_{k\text{ slots for covectors}},
			\underbrace{
				\underline{\phantom{x}},
				\underline{\phantom{x}},
				\underline{\phantom{x}},
				\ldots,
				\underline{\phantom{x}},
				\underline{\phantom{x}},
				\underline{\phantom{x}}
			}_{l\text{ slots for vectors}},
			)
		\]
		Think of it as a machine, that takes in $k$ covectors, and $l$
		vectors, and spits out a real number. 

		\paragraph{Multilinearity} A tensor has to be multilinear, meaning that the tensor is linear in
		each slot separately. For example, if the first covector is varied,
		\begin{align*}
			T(\alpha\underline{\omega}+\beta\underline{\xi},
			\underline{\omega}_2,\ldots,\underline{\omega}_k,
			\overline x_1,\ldots,\overline x_l)
			&=
			\alpha
			T(\underline{\omega},\underline{\omega}_2,\ldots,
			\underline{\omega}_k,\overline x_1,\ldots,\overline x_l)
			\\
			&\quad+
			\beta
			T(\underline{\xi},\underline{\omega}_2,\ldots,
			\underline{\omega}_k,\overline x_1,\ldots,\overline x_l).
		\end{align*}
		
		In general, for $\omega^{(\mu)}\in V^*$ and $v^{(\mu)} \in V$, we can express it as (explicitly putting sums here)
		\begin{align*}
			&T\left(
			\sum_a \alpha_a \omega_a^{(1)},
			\ldots,
			\sum_a \alpha_a \omega_a^{(k)},
			\sum_b \beta_b v_b^{(1)},
			\ldots,
			\sum_b \beta_b v_b^{(l)}
			\right)
			=\\
			&\sum_{a_1,\ldots,a_k}
			\sum_{b_1,\ldots,b_l}
			\alpha_{a_1}\cdots\alpha_{a_k}
			\beta_{b_1}\cdots\beta_{b_l}
			T\left(
			\omega_{a_1}^{(1)},\ldots,\omega_{a_k}^{(k)},
			v_{b_1}^{(1)},\ldots,v_{b_l}^{(l)}
			\right).
		\end{align*}

	\subsubsection{Vectors and covectors as tensors}
		From the definition, immediately we can see that a covector
		is a tensor.
		\[
			\underline{\omega} \in V^*,\qquad \underline{\omega}:V\to\R
		\]
		Specifically, it is a type $\binom{0}{1}$ tensor\footnote{$(k,l)$ and $\binom{k}{l}$ are interchangeable}. It takes one vector, 
		and spits out a real number. We refer to it as belonging to the type (0,1) tensor. 
		\[
			\underline\omega \in \binom{0}{1}
		\]

		Using isomorphism, and the properties and identities shown before, the vector space can
		be thought of as \textit{accepting} dual vectors, because the dual space of the dual space
		is a vector space. So it has one slot for a covector. 
		\[
			\overline{\lambda}\in V,\qquad \overline{\lambda}: V^* \to \R
		\]
		That is, a vector is an element of type $\binom{1}{0}$ tensors.
		\[
			\overline{\lambda}\in\binom{1}{0}
		\]
		
		From that logic, a scalar has no slots for a vector or covector, and simply returns a real number.
		Therefore, a scalar is a $\binom{0}{0}$ tensor.
		\[
			a\in\R, \qquad a:\R\to\R \implies \quad a\in \binom{0}{0}
		\]

	\subsubsection{Tensors as vector/covector operators}
		Suppose a tensor $T\in\binom{1}{1}$, such that
		\[
			T: V^*\times V\to \R
		\]
		Note that because it is multilinear, we can sort of 
		leave a slot \textit{blank} to be filled in. We can
		fill the first slot with a covector, and leave the vector 
		slot as blank.
		\[
			T(\underline\omega,\underline{\phantom{x}})
		\]
		It should be immediately obvious then, that this tensor, left with a vacant slot, becomes
		an independent operator. It becomes a dual vector, since it accepts only one vector.
		\[
			T(\underline\omega,\underline{\phantom{x}})
			\in V^*
		\]
		That is, it is a map, from covectors to covectors, when the vector slot is left vacant.
		It accepts a covector, and turns it into another covector.
		\[
			T:V^* \to V^*,\qquad 
			T(\underline\omega,\underline{\phantom{x}})
			\in\binom{0}{1}
		\]
		Similarly, if the same tensor were supplied with only a vector, it has one remaining slot
		for a covector, and hence, it becomes a vector. It takes a vector, and returns another vector.
		\[
			T(
				\underline{\phantom{\omega}},
				\overline{v}
			)
			\in V,\quad T:V \to V ,\qquad 
			T(
				\underline{\phantom{\omega}},
				\overline{v}
			)
			\in\binom{1}{0}
		\]
		In general, for a $\binom{k}{l}$ tensor, when supplied with $m$ covectors and $n$ vectors,
		becomes an type $(k-m, l-n)$ tensor, with condition $m\leq k$ and $n\leq l$. 
		\[
			T\in\binom{k-m}{l-n}
		\]

		The point here is filling slots of tensors 
		produces a tensor of lower type.

		\paragraph{Tensor Rank} The rank or order of a tensor is the total number of slots it has
		for accepting other tensors. For a $(k,l)$ tensor, the rank is
		\[
			T\in\binom{k}{l},\qquad \rank T = k+l
		\]
		For example, the Riemman tensor $\tensor{R}{^\rho_{\sigma\mu\nu}}$, of type (1,3), is a rank 4 tensor. The Levi-Civita symbol $\tensor{\varepsilon}{_{ij}^k}$ is a rank 3 tensor, since it has 3 total slots.
		The Kronecker-delta $\tensor{\delta}{^i_j}$, metric tensor $g_{\mu\nu}$, and matrices $a_{ij}$ in general are rank 2 tensors, because
		each has two tensor slots. Vectors, $v = v^\mu \p_\mu$ and duals, $\omega = \omega_\mu dx^\mu$, are rank 1 tensors, because
		it can only accept one tensor, and finally, a scalar is a
		rank 0 tensor, because it accepts no tensors. 

		\paragraph{Examples}
		\begin{enumerate}[(1)]
			\item An $n\times n$ matrix.
				
				All square matrices are isomorphic to type $(1,1)$ tensors. They
				represent a tensor of type (1,1), once bases for $V$ and $V^*$ are chosen. Note 
				that they themselves are not tensors in the geometric sense, just as row and column matrices
				are not covectors and vectors, but they are isomorphic to each other, and
				thus we can perform the same operations on them.

				From the previous chapter, we discussed how column matrices 
				are isomorphic to vectors.
				\[
					\bmqty{
						v^1\\v^2\\\vdots\\v^n
					}\in V
				\]
				and row matrices are isomorphic to covectors.
				\[
					\bmqty{
						\omega_1 &
						\omega_2 &
						\cdots &
						\omega_n
					}
					\in V^*
				\]
				Now consider how do we make a matrix output a real number? Let's do a 
				$2\times 2$ matrix to keep the math simple.

				If you feed a $2\times 2$ matrix a column vector, it becomes a column
				vector
				\[
					\bmqty{
						a_{11} & a_{12}\\
						a_{21} & a_{22}\\
					} 
					\bmqty{
						v_1\\ v_2	
					} 
					=
					\bmqty{
						v_1 a_{11} + v_2 a_{12}\\
						v_1 a_{21} + v_2 a_{22}\\
					} 
				\]
				This is exactly analogous to what we have shown earlier. A (1,1) tensor,
				fed a vector, becomes a vector. This tensor is then a map from a vector, to
				another vector. If you then operate it to a covector, it returns a real number.

				For a second example, consider multiplying a row matrix instead. Doing 
				so returns another row matrix.
				\[
					\bmqty{
						\omega_1& \omega_2	
					} 
					\bmqty{
						a_{11} & a_{12}\\
						a_{21} & a_{22}\\
					} 
					=
					\bmqty{
						\omega_1 a_{11} +\omega_2 a_{21}&
						\omega_1 a_{12} +\omega_2 a_{22}
					} 
				\]
				Again, this is analogous to what we have shown. A tensor, when fed 
				a covector, returns another covector. It maps a covector to another 
				covector. If you operate this covector to a vector, it returns a real number. 

				Now consider feeding the matrix both a vector and a covector. You will find that
				\[
					\bmqty{
						\omega_1& \omega_2	
					} 
					\bmqty{
						a_{11} & a_{12}\\
						a_{21} & a_{22}\\
					} 
					\bmqty{
						v_1\\ v_2	
					} 
					=
						v_1 \omega_1 a_{11} + v_1 \omega_2 a_{21} +
						v_2 \omega_1 a_{12} + v_2 \omega_2 a_{22}
				\]

				Therefore, it returned a scalar, which is precisely what a tensor is defined to do.
				It accepted one covector, and one vector, and returned a real number.
			\item A cross product

				Consider the introductory Mathematical Physics definition of a cross product.
				\[
					\vb{u}\times \vb{v} = \varepsilon_{ijk} u_i v_j \vu{e}_k
				\]
				Now we use the established tensor notation, of raised and lowered indices
				indicating vector and covector components. We have vectors $u = u^i\p_i$, $v = v^j\p_j$,
				and the resulting cross product to be $(u\times v)^k \p_k$. Thus, calculating the cross
				product components\footnote{Shown on the next page}, we have
				\[
					(u\times v)^k = 
					\tensor{\varepsilon}{_{ij}^k} u^i v^j
				\]
				
				We have shown then, that the cross product of two vectors
				returns a vector. But that all hinges on the fact that the 
				Levi-Civita symbol is a tensor\footnote{Technically it's a 
				pseudotensor which is only really a tensor under orthogonal
				transformations and the Euclidean isometry group, i.e.,
				in flat, Euclidean metric and cartesian coordinates. It is a
				tensor \textit{density} of weight $+1$, and under general coordinate
				transformation, it changes, but since this is an example, we'll 
				call it a tensor for now}. The 
				Levi-Civita tensor is a rank 3 tensor, of type (1,2).
				It accepts 
				two vectors and one covector. Once it gets fed
				two vectors, its vector slots are filled, and thus it can only accept 
				a covector. The two lower slots are contracted with the vectors $u$ and $v$, leaving the upper $k$ index free.
				The resulting object is therefore a linear map from covectors to real numbers, which is naturally
				an element of $V^{**}\cong V$, i.e., a vector. 

				Therefore, it then acts as a map from two vectors to another vector,
				when the two vectors are supplied, and returns a vector itself. This is
				precisely the cross product.

				Note that since it is a rank 3 tensor, and it ate two rank 1 tensors, its rank
				also went to 1. When a rank $n$ tensor accepts $m$ vectors/covectors, its rank
				decreases by $n-m$.
		\end{enumerate}

		\begin{remark}
			An explicit cross product is expressed as follows. Note that this
			only applies in the Euclidean metric. For two vectors
			$u = u^i \p_i$ and $v = v^j \p_j$, we aim to find the components
			of $u \times v$. Recall that the cross product of the basis
			vectors is
			\[
				\p_i \times \p_j = \tensor{\varepsilon}{_{ij}^k} \p_k
			\]
			Then
			\begin{align*}
				u \times v
				&= (u^i \p_i) \times (v^j \p_j) \\
				&= u^i v^j (\p_i \times \p_j) \\
				&= u^i v^j \tensor{\varepsilon}{_{ij}^k} \p_k \\
				&= \tensor{\varepsilon}{_{ij}^k} u^i v^j \p_k \\
				&= (u \times v)^k \p_k
			\end{align*}
			Therefore the components are
			\[
				(u \times v)^k = \tensor{\varepsilon}{_{ij}^k} u^i v^j
			\]
		\end{remark}

		\paragraph{Ordering of Covectors and Vectors}
		An important property is the following: the \emph{order} of the covector
		and vector slots matters. One may reorder the slots as long as the
		relative order among the covectors and among the vectors is preserved,
		but changing the intrinsic order \emph{within} either group is not allowed.

		For example, consider a $(3,3)$ tensor
		\[
			T(\omega_1, \omega_2, \omega_3, v^1, v^2, v^3).
		\]
		One may interleave the vectors and covectors, since the two groups
		occupy independent sets of slots, e.g.,
		\[
			T(\omega_1, \omega_2, \omega_3, v^1, v^2, v^3)
			=
			T(\omega_1, \omega_2, v^1, \omega_3, v^2, v^3)
			=
			T(\omega_1, \omega_2, v^1,v^2,v^3, \omega_3)
		\]
		However, permuting the covectors among themselves, or the vectors among
		themselves, changes the tensor's action:
		\[
			T(\omega_1, \omega_2, \omega_3, v^1, v^2, v^3)
			\neq
			T(\omega_1, \omega_2, \omega_3, v^2, v^1, v^3)
			\;\neq\;
			T(\omega_1, \omega_2, v^3, \omega_3, v^1, v^2).
		\]

% =======================================
% MEETING 9
% ---------------------------------------
% DATE   - September 18, 2026
% FORMAT - F2F
% =======================================
% TOPIC  - Tensor product
% =======================================
\subsubsection{Tensor Product}
	\label{subsec:tensor product}
	The point of a tensor is this: A tensor is a function (or a \textit{machine}, if you will),
	that accepts covectors and vectors with $k$ and $l$ slots. After all the slots are filled,
	then the tensor spits out a real number. However, as long as those slots aren't filled yet,
	it becomes a tensor of lower rank, or lower type.

	Ordering matters between the sets of vectors and covectors. You
	can permute between them, but permuting within the group is not allowed. 

	Recall the identity we set out earlier. We have the following 
	coordinate basis for covectors and vectors:
	\begin{align*}
		\overline{e}_\mu = \p_\mu && \underline{e}^\mu = dx^\mu
	\end{align*}

	We can find the components of an arbitrary vector or covector likeso:
	\[
		v^\mu = dx^\mu[v] \equiv v[dx^\mu]
	\]
	and
	\[
		\omega_\mu = \omega[\p_\mu]\equiv \p_\mu[\omega]
	\]
	This is a sort of product of operation which is symmetric, hence it is called the 
	interior product.
	Coacting, or \textit{multiplying} rank 1 tensors together
	results on another tensor (recall: $a\in\R$ is just a rank 0 tensor).

	We need to define the notion of a tensor product, which
	allows us to combine tensors of different rank to produce a 
	new tensor. 

	\begin{definition}{Tensor Product}
	    If $\varphi$ and $\vartheta$ are covectors, or $(0,1)$ tensors, then
		their tensor product $\varphi\otimes\vartheta$ is a $(0,2)$ tensor, such that, for $x,y\in(1,0)$:
		\begin{equation}
			\label{eq:tensor-product-defn}
			(\varphi\otimes\vartheta)[x,y] = \varphi[x]\vartheta[y]
		\end{equation}

		More generally, if $T$ has type $(k, l)$ and $S$ has type $(m, n)$,
		then $T \otimes S$ has type $(k + m, l + n)$, defined by
		\begin{align*}
			&(T \otimes S)(\omega_1, \ldots, \omega_k, \omega'_1, \ldots, \omega'_m,
			v_1, \ldots, v_l, v'_1, \ldots, v'_n)
			=\\ &T(\omega_1, \ldots, \omega_k, v_1, \ldots, v_l)
			  \,S(\omega'_1, \ldots, \omega'_m, v'_1, \ldots, v'_n).
		\end{align*}
		This is also known as an outer product.
	\end{definition}

	\paragraph{Example}

	Say we have one-forms $dx^1$ and $dx^2$. If we let them act on a pair 
	of vetors $x$ and $y$, we have
	\begin{align*}
		(dx^1\otimes dx^2) [x,y] &= \dd{x}^1[x]\dd{x^2}[y] \\
		&=x^1 y^2
	\end{align*}
	Those are individual components.
	
	It should be clear from the definition that the outer product is
	not commutative.
	\[
		\varphi\otimes\vartheta 
		\neq 
		\vartheta\otimes\varphi
	\]
	For example, if we let it act on the vectors $x,y$, then
	\begin{align*}
		\varphi\otimes\vartheta[x,y] &=
		\varphi[x]\vartheta[y]\\
		\vartheta\otimes\varphi[x,y]&=
		\vartheta[x]\varphi[y]\\
		\varphi[x]\vartheta[y]&\neq
		\vartheta[x]\varphi[y]
	\end{align*}
	\paragraph{Properties} The tensor product must satisfy the following 
	properties
	\begin{enumerate}[(1)]
		\item \text{Bilinearity}

			For scalars $a, b$ and tensors
			  $\varphi, \theta, \omega$ of compatible type,
			  \[
				  \varphi \otimes (a\theta + b\omega)
				  = a\,(\varphi \otimes \theta) + b\,(\varphi \otimes \omega),
			  \]
			  \[
				  (a\varphi + b\theta) \otimes \omega
				  = a\,(\varphi \otimes \omega) + b\,(\theta \otimes \omega).
			  \]
		\item \text{Associativity}

			\[(\varphi \otimes \theta) \otimes \omega
			= \varphi \otimes (\theta \otimes \omega)\]
		\item \text{Multiplicative scaling}

			\[(a\varphi) \otimes \theta
			   = a\,(\varphi \otimes \theta)
		   = \varphi \otimes (a\theta)\]
		\item \text{Noncommutativity}
			
			In general, for tensors $\phi$ and $\theta$ acting on arbitary tensors,
			\[\varphi \otimes \theta \neq \theta \otimes \varphi\]
	\end{enumerate}
	\begin{proof}
	   Exercise to the student. 
	\end{proof}

\subsubsection{Basis Tensors}
	Like vectors and covectors, tensors which are defined as a linear combination of 
	scalars by basis vectors or basis covectors, we must find a way
	to generalize what a basis is for higher order tensors.

	Consider a tensor $g \in \binom{0}{2}$, acting on two vectors
	$x, y \in \binom{1}{0}$, such that
	\[
		g[x, y] \mapsto \alpha \in \R.
	\]
	We can express $x$ and $y$ in terms of a basis:
	\[
		x = x^\mu \p_\mu,
		\qquad
		y = y^\nu \p_\nu.
	\]
	Thus, we can express $g$ as
	\[
		g[x, y]
		= g\bqty{x^\mu \p_\mu,\; y^\nu \p_\nu}.
	\]
	Since $x^\mu$ and $y^\nu$ are just numbers, i.e., constants, we
	can pull them out:
	\[
		g[x, y]
		= x^\mu y^\nu\, g[\p_\mu, \p_\nu].
	\]
	Notice that $g$ is accepting two vectors. Hence, the entire
	expression is just a product of real numbers. Let us define
	\[
		g[\p_\mu, \p_\nu] = g_{\mu\nu}.
	\]
	Thus we can rearrange it to
	\[
		g[x, y]
		= \tensor{g}{_{\mu\nu}}\, x^\mu y^\nu.
	\]
	Recall that $x^\mu = dx^\mu[x]$ and $y^\nu = dx^\nu[y]$. Then
	\[
		g[x, y]
		= \tensor{g}{_{\mu\nu}}\,
		  dx^\mu[x]\, dx^\nu[y].
	\]
	From the definition of the tensor product,
	\[
		dx^\mu[x]\, dx^\nu[y]
		= (dx^\mu \otimes dx^\nu)[x, y].
	\]
	We can then represent the entire expression as
	\[
		g[x, y]
		= \tensor{g}{_{\mu\nu}}\,
		  (dx^\mu \otimes dx^\nu)[x, y].
	\]
	However, the vectors $x$ and $y$ are arbitrary. We can remove
	them, hence
	\[
		g = g_{\mu\nu}\, dx^\mu \otimes dx^\nu.
	\]
	Notice here that we have represented a tensor
	$g \in \binom{0}{2}$ as a linear combination of
	$\binom{0}{2}$ tensors $dx^\mu \otimes dx^\nu$, with
	components $g_{\mu\nu}$.

	Recall what we have done in the previous sections. A vector, for
	instance, can be represented by a product of components and a
	basis, and similarly for covectors:
	\[
		v = v^\mu \p_\mu,
		\qquad
		u = u_\nu\, dx^\nu.
	\]
	We have arrived at a similar expression. Therefore, we can
	conclude that $dx^\mu \otimes dx^\nu$ must be a basis. This is
	what we call a \textit{basis tensor}.

	In the same logic, if you have a tensor $m \in \binom{2}{0}$
	acting on two covectors in $\binom{0}{1}$, then we arrive at the
	same conclusion:
	\[
		m = m^{\mu\nu}\, \p_\mu \otimes \p_\nu,
	\]
	where $m^{\mu\nu} = m[dx^\mu, dx^\nu]$. Therefore
	$\p_\mu \otimes \p_\nu$ is also a basis tensor.

	\begin{definition}{Basis Tensor Representation}
		In general, we can represent an arbitrary $(k, l)$ tensor
		$T$ as
		\[
			T
			=
			\tensor{T}{^{a_1 a_2 \cdots a_k}_{b_1 b_2 \cdots b_l}}
			\p_{a_1} \otimes \p_{a_2} \otimes \cdots \otimes \p_{a_k}
			\otimes
			dx^{b_1} \otimes dx^{b_2} \otimes \cdots \otimes dx^{b_l}.
		\]
		where
		\[
			\p_{a_1} \otimes \cdots \otimes \p_{a_k}
			\otimes
			dx^{b_1} \otimes \cdots \otimes dx^{b_l}.
		\]
		is a $(k, l)$ basis tensor, and
		\[
			\tensor{T}{^{a_1 \cdots a_k}_{b_1 \cdots b_l}}
		\]
		are the $(k, l)$ tensor components.
	\end{definition}

	As with vectors and covectors, basis tensors must follow the
	property
	\[
		e^\mu[e_\nu] = e_\nu[e^\mu] = \tensor{\delta}{^\mu_\nu}.
	\]
	An equivalent statement for a general basis tensor is that the
	$(k, l)$ basis tensor acts on the corresponding basis vectors and
	covectors to give products of Kronecker deltas:
	\begin{align*}
		&\pqty{
			\p_{a_1} \otimes \cdots \otimes \p_{a_l}
		      \otimes 
			dx^{b_1} \otimes \cdots \otimes dx^{b_k}
		}
		\bqty{
			dx^{c_1}, \ldots, dx^{c_l}
			,\;
			p_{d_1}, \ldots, \p_{d_k}
		}\\
		&=
		\tensor{\delta}{^{c_1}_{a_1}} \cdots \tensor{\delta}{^{c_k}_{a_k}}\,
		\tensor{\delta}{^{b_1}_{d_1}} \cdots \tensor{\delta}{^{b_l}_{d_l}}.
	\end{align*}
	This is the natural generalization of $e^\mu[e_\nu] =
	\delta^\mu{}_\nu$ to higher-rank basis tensors: the covector
	slots and vector slots each pick off one Kronecker delta, and
	the full result is the product. It is this property that makes
	the representation
	\[
		T
		=
		\tensor{T}{^{a_1 \cdots a_k}_{b_1 \cdots b_l}}
		\,
		\p_{a_1} \otimes \cdots \otimes \p_{a_k}
		\otimes
		dx^{b_1} \otimes \cdots \otimes dx^{b_l}
	\]
	unique, and it is what allows the components to be extracted by
	the rule
	\[
		\tensor{T}{^{a_1 \cdots a_k}_{b_1 \cdots b_l}}
		=
		T\bqty{dx^{a_1}, \ldots, dx^{a_k},\;
		       \p_{b_1}, \ldots, \p_{b_l}}.
	\]
	In other words, the $(a_1, \ldots, a_k;\, b_1, \ldots, b_l)$ component
	of $T$ is obtained by feeding $T$ the corresponding basis
	covectors and basis vectors.
	
	This sort of notation is relatively uncommon in mathematics, but has been 
	universally accepted in physics.

	\paragraph{Example} Consider a $(2,1)$ tensor represented in tensor basis form
	\[
		\kappa = \tensor{\kappa}{^{\mu\nu}_{\sigma}} \p_\mu \otimes \p_\nu \otimes dx^\sigma
	\]
	If we want to find the $\tensor{T}{^{\alpha\beta}_\gamma}$-th component, we feed it the 
	basis tensors $dx^\alpha, dx^\beta,\text{ and }\p_\gamma$. 
	\begin{align*}
		\kappa
		[dx^\alpha,dx^\beta, \p_\gamma] 
		&= 
		\tensor{\kappa}{^{\mu\nu}_{\sigma}} 
		\p_\mu \otimes \p_\nu \otimes dx^\sigma
		[dx^\alpha,dx^\beta, \p_\gamma] \\
		&=
		\tensor{\kappa}{^{\mu\nu}_{\sigma}} 
		\p_\mu[dx^\alpha]
		\p_\nu[dx^\beta]
		dx^\sigma[\p_\gamma]\\
		&=
		\tensor{\kappa}{^{\mu\nu}_{\sigma}} 
		\tensor{\delta}{^\alpha_\mu}
		\tensor{\delta}{^\beta_\nu}
		\tensor{\delta}{^\sigma_\gamma}\\
		&=
		\tensor{\kappa}{^{\alpha\nu}_{\sigma}} 
		\tensor{\delta}{^\beta_\nu}
		\tensor{\delta}{^\sigma_\gamma}\\
		&=
		\tensor{\kappa}{^{\alpha\beta}_{\sigma}} 
		\tensor{\delta}{^\sigma_\gamma}
		\\
		&=
		\tensor{\kappa}{^{\alpha\beta}_{\gamma}} 
	\end{align*}
	
	Thus, we can see that the same contraction operation applies here as it 
	does for vectors and covectors.
\subsubsection{Symmetric Tensor Product and Transformation Law}
	% We mentioned several lectures ago that 
	% to get the notion of lengths and orthogonality 
	% between vectors, we need an operation that will
	% allow us to measure lengths and angles. An operation,
	% which takes two vectors, and spits out a scalar. 
	% We shall discuss that operation now.

	A tensor product, in general, is noncommutative.
	\[
		\p_\mu \otimes \p_\nu \neq \p_\nu \otimes \p_\mu
	\]
	However, there are important tensors for which interchanging the
	two tensor factors leaves the tensor unchanged. Such tensors are
	called \textit{symmetric tensors}. These tensors have a notion of commutativity.	

	Consider a tensor $g\in(2,0)$. 
	\[
		g = 
		g_{\mu\nu} dx^\mu \otimes dx^\nu
	\]
	In a 2D manifold (to keep the math simple), its explicit decomposition
	would be
	\begin{align*}
		g = 
		g_{11} \dd{x}^1 \otimes dx^1 +
		g_{12} \dd{x}^1 \otimes dx^2 +
		g_{21} \dd{x}^2 \otimes dx^1 +
		g_{22} \dd{x}^2 \otimes dx^2
	\end{align*}
	Now, suppose the tensor is symmetric, i.e.,
	\[
		g_{\alpha\beta} = g_{\beta\alpha} 
	\]
	This is the same sense of a matrix being symmetric. There
	are many cases where the tensor you are working with are symmetric,
	for example, in GR, the general metric tensor is symmetric. The 
	stress tensor $\sigma_{ij}$ and strain tensor $\varepsilon_{ij}$ is symmetric, the inertia tensor $I$, 
	and so on.

	In that case, we can write the explicit decomposition as
	\[
		g = 
		g_{11} \dd{x}^1 \otimes dx^1 +
		g_{12} (dd{x}^1 \otimes dx^2 +
		\dd{x}^2 \otimes dx^1) +
		g_{22} \dd{x}^2 \otimes dx^2
	\]
	Notice, however, that all the tensor elements are now symmetric ($g_{11}=g_{11}, g_{22}=g_{22}$, and $g_{12}=g_{21}$, by construction).

	We have a sense of a symmetric product here, which we can use. It suggests that,
	for symmetric tensors, it is convenient to suppress the explicit tensor product
	and introduce an implicit symmetric product. Let us define the following

	\begin{definition}{Symmetric Tensor Product}
		For a symmetric tensor, i.e.,
		\[
			dx^{\mu}\otimes dx^\nu = dx^\nu \otimes dx^\mu
		\]
		We define a symmetric tensor product $dx^\mu dx^\nu$ such that
		\begin{align*}
			dx^{\mu}dx^{\nu} &= 
			\frac{1}{2}\pqty{
				dx^{\mu}\otimes dx^\nu + dx^\nu \otimes dx^\mu
			}\\
							 &= dx^\nu dx^\mu
		\end{align*} 
		This is also known as the symmetrized tensor product.
	\end{definition}

	Now that we can do that, we can rewrite $g$ as 
	\[
		g = g_{11}(dx^1)^2 + 2g_{12} dx^1 dx^2 + g_{22} (dx^2)^2
	\]
	Notice, if you interpret this physically, that 
	this is just an infinitesimal square of arclength.

	Now that we can write it in an easier geometric way, we just collapse the expression back to
	the general tensor product, with the understanding that the product of the differentials
	represents an appropriate symmetric tensor product.
	\[
		g = g_{\mu\nu} dx^\mu dx^\nu
	\]
	Notice that this is just a normal covector operation. We can examine
	how the tensor components $g_{\mu\nu}$ transform by performing a change 
	of basis, from $x^\mu$ to $y^\sigma$. 

	Let $x^\mu = x^\mu(y^\sigma)$. Recall the change of basis formulation defined in the last chapter. That is,
	\[
		dx^\mu = \pdv{x^\mu}{y^\sigma} dy^\sigma
	\]
	We can then write the expression as
	\begin{align*}
		g 
		&= g_{\mu\nu} dx^\mu dx^\nu\\
		&= g_{\mu\nu} 
		\pqty{
			\pdv{x^\mu}{y^\sigma} 
			dy^\sigma 
		}
		\pqty{
			\pdv{x^\nu}{y^\rho} 
			dy^\rho 
		}
	\end{align*}
	The Jacobians are just constants, we can pull them out
	\[
		g = 
			\pdv{x^\mu}{y^\sigma} 
			\pdv{x^\nu}{y^\rho} 
			\,
		g_{\mu\nu} 
			dy^\sigma
			dy^\rho 
	\]
	On the other hand, since $g$ is the same tensor regardless of the 
	coordinate system, we may express it directly in terms of the new coordinates
	as
	\[
		g = g'_{\sigma\rho} dy^\sigma dy^\rho
	\]
	Comparing the two expressions, we now have a tensor transformation rule. This
	is how the components of a $(0,2)$ tensor transforms under a change of basis.
	\[
		g'_{\sigma\rho} 
		=
		\pdv{x^\mu}{y^\sigma} 
		\pdv{x^\nu}{y^\rho} 
		\,
		g_{\mu\nu} 
	\]
	This is the transformation law for a $(0,2)$ symmetric tensor. In particular,
	the metric components transform covariantly because each lower index
	contributes one factor of the Jacobian.

	This all hinges on the assumption that the tensor product is symmetric, and the tensor components are symmetric.
	The tensor product is hence, implicit.

\subsubsection{Contraction Operation}
	We need to define what a contraction is.
	
	\begin{definition}{Contraction}
		If a tensor $T$ has a contravariant index and a covariant index
		which are contracted, then the contraction operation $C$ acts on
		the tensor by summing over the repeated index. If
		$T\in{(k,l)}$, then
		\[
			C(T)\in{(k-1,l-1)}.
		\]
		Thus, $CT$ is a new tensor with both its contravariant and
		covariant rank reduced by $1$.
	\end{definition}
	For example, suppose $T$ has components
	\[
		T
		=
		T^{a_1\cdots \mu\cdots a_k}{}_{b_1\cdots \mu\cdots b_l}
		\,
		\p_{a_1}\otimes\cdots\otimes \p_\mu\otimes \cdots \otimes\p_{a_k}
		\otimes
		dx^{b_1}\otimes\cdots\otimes dx^\mu\otimes\cdots\otimes dx^{b_l},
	\]
	where the repeated index $\mu$ appears once upstairs and once
	downstairs. The contraction sums over all possible values of
	$\mu$:
	\[
		CT
		=
		\sum_{\mu}
		T^{a_1\cdots \mu\cdots a_k}{}_{b_1\cdots \mu\cdots b_l}
		\,
		\p_{a_1}\otimes\cdots
		\otimes
		dx^{b_1}\otimes\cdots.
	\]
	The contracted index $\mu$ then disappears, leaving a tensor of
	type $(k-1,l-1)$.

	Another example, consider a tensor $T\in(1,1)$ with basis
	$dx^\mu$ and $\p_\mu$:
	\[
		T = T^\mu{}_\nu\,\p_\mu\otimes dx^\nu.
	\]
	When we perform a contraction, the vector and covector basis
	are evaluated against one another:
	\[
		dx^\nu[\p_\mu]=\delta^\nu_\mu.
	\]
	Thus,
	\[
		CT
		=
		T^\mu{}_\nu\,\delta^\nu_\mu
		=
		T^\mu{}_\mu,
	\]
	which is a scalar, of type $(0,0)$. This is the trace of the $(1,1)$ tensor.
	
\subsubsection{Abstract Index Notation}
	When it comes to tensor notation, there are several competing schools of thought
	regarding how it would be represented. We can show these 
	different notation with the following
	\begin{enumerate}[(1)]
		\item Index--free notation

			To illustrate this, let's show the following. Consider a tensor $S\in\binom{p}{q}$, and 
			another tensor $T\in\binom{r}{s}$. 

			Individually, the tensors act like so
			\[
				S[\theta_1,\cdots,\theta_p,\phi_1,\cdots,\phi_q]
			\]
			and
			\[
				T[\theta_1,\cdots,\theta_r,\phi_1,\cdots,\phi_s]
			\]
			The tensor product
			between them is
			\[
				S\otimes T \in \binom{p+r}{q+s}
			\]
			That is, it accepts $p+r$ covectors, and $q+s$ vectors. Let $\theta_i$ be covectors and 
			$\phi_j$ be vectors. 
			Then, the tensor product acts like
			\[
				(S\otimes T)[\theta_1,\cdots, \theta_{p},\phi_1,\cdots,\phi_q,\theta_{p+1},\cdots,\theta_{p+r},\phi_{q+1},\cdots,\phi_{q+s}]
			\]
			This can all be confusing, no? It is easy to make mistakes in this notation. As such, 
			you need to respect the order of the covectors and vectors. 
			We can rearrange the order by which the tensor accepts arguments. This is simply equivalent to
			\[
				(S\otimes T)[\theta_1,\cdots, \theta_{p+r}, \phi_1, \cdots, \phi_{q+s}]
			\]
			Because the order of the duals and vectors themselves do not matter, what matters is 
			the order of the vectors among the group of vectors, and the order of the covectors among the 
			group of covectors. That order is respected, but to simplify things, the new tensor acts the same way
			if we feed it the $S$ covectors first, then the $S$ vectors, followed by the $T$ covectors and vectors. 

			Writing it down in component form, we have
			\begin{align*}
				&S\otimes T =
				\tensor{S}{^{i_1,\cdots,i_p}_{j_1,\dots,j_q}}
				\tensor{T}{^{k_1,\cdots,k_q}_{l_1,\dots,l_s}}\\
						   &\p_{i_1}\otimes\cdots\otimes\p_{i_p}\otimes dx^{j_1}\otimes\cdots\otimes dx^{j_q}
				\p_{k_1}\otimes\cdots\otimes\p_{k_q}\otimes dx^{l_1}\otimes\cdots\otimes dx^{l_s}
			\end{align*}
			Equivalently,
			\begin{align*}
				&S\otimes T =
				\tensor{S}{^{i_1,\cdots,i_p}_{j_1,\dots,j_q}}
				\tensor{T}{^{k_1,\cdots,k_q}_{l_1,\dots,l_s}}\\
						   &\p_{i_1}\otimes\cdots\p_{i_p}\otimes
				\p_{k_1}\otimes\cdots\p_{k_q}\otimes
				dx^{j_1}\otimes\cdots\otimes dx^{j_q}\otimes
				dx^{l_1}\otimes\cdots\otimes dx^{l_s}
			\end{align*}
			We do this because the tensor product commutes for vectors and covectors.
			The slots for covectors and the slots for vectors should not be confused with 
			one another. Covector order and vector order matters, but the order of the covectors and vectors themselves don't really matter.

			Under this notation, there is no ambiguity. As long as the order of the covectors and vectors 
			are respected, we can swap them and indicate which slots belong to which.

			However, as you might have noticed, this notation is very cumbersome. I am intentionally
			showing you how inconvenient it is to use this notation. All those shit regarding 
			bases that you need to follow. However, mathematicians love this notation.
			They prefer it precisely because it is rigourous, and no confusion can happen.
			
			This notation is called the \textit{index-free} notation. It is index free because the 
			tensor itself has no index, only the component has. It is what we have been
			using so far to understand tensors. However, it has a few drawbacks.

			First, is that the expression of a tensor is very long and very cumbersome to write. It is
			easy to make mistakes. Second, is that you do not know the type of the tensor immediately just
			by looking at it. We need to specify that a tensor is of type $(k,l)$, which is separate
			from how we define the tensor itself. It is difficult to perform calculus on it, and 
			it's just inconvenient. Contraction operations on tensors using this notation is hard.
			Manipulating things is diffilcult. For example, showing that
			\[
				S\otimes T \neq T\otimes S
			\]
			explicity is cumbersome, since we need to define everything.

			However, this is very common in pure mathematics. It is inconvenient, but rigourous. It is unambigous,
			and it is exact.
		\item Index notation

			This notation is quite old fashioned. It is used by 
			Physicists to express and  manipulate tensors, specially from those
			under the old universities such as Princeton or MIT. It was widely
			accepted in physics in the last century. Here, we denote the tensor with indices.

			For example, a $(2,1)$ tensor is just 
			\[
				\tensor{T}{^{\mu\nu}_{\sigma}}
			\]

			That's the tensor itself. All the basis that accompany it, and its components, are implied. 

			A tensor product between, say $\tensor{S}{^\mu_\nu}$ and $\tensor{T}{^{\alpha\beta^\gamma}_{\epsilon\xi}}$
			is just
			\[
				\tensor{S}{^\mu_\nu}
				\otimes
				\tensor{T}{^{\alpha\beta\gamma}_{\epsilon\xi}}
				=
				\tensor{(S\otimes T)}{^{\mu\alpha\beta\gamma}_{\nu\epsilon\xi}}
			\]

			It has a few advantages. First, you already know the type of the tensor without being explicitly
			told. From the symbol itself, for example, we can already tell that the tensor above is a $(4,3)$ tensor.
			Second, manipulating it for calculus, contraction, and other operations are easy. There is no need
			to fuss about with the basis vectors or covectors. It is convenient.

			However, it also has some drawbacks. You can't easily tell if what you are manipulating 
			is the tensor itself, or just its components. It's hard to check the validity of the mathematics
			you do if it is difficult to identify what is it you are working with. Also, if you 
			manipulate larger tensors, the amount of indices become too many, and it might be 
			more difficult to work with. Lastly,
			mathematicians will not understand what you are referring to. If you give a mathametician, say
			\[
				\tensor{T}{^\mu_\beta}
			\]
			They will not call that a tensor. They will call that a scalar. It's the operation of a tensor
			on the bases. It is a collection of components.
			\[
				T[dx^\beta,\p_\mu]
			\]
			So for them, what physicists refer to as a tensor is not necesarilly a tensor itself. 
			So there is some difficulty collaborating between the two.

		\item Abstract Index notation

			Penrose and others, specially from those that came from the Chicago school of thought,
			went halfway between the two. They used both elements of the notation.

			They \textit{defined} the following: If a tensor symbol uses \textit{Latin} indices,
			it is a tensor, the geometric object.
			
			For example,
			\[
				\tensor{T}{^{a}_{bc}}
			\]
			From this, it is immediately clear that the tensor is a $(1,2)$ tensor. There is no
			explicit reference to its type, because it is obvious from the indices alone.

			A tensor product is simply
			\[
				\tensor{T}{^{ab}_{c}}
				\tensor{S}{^d_e} 
				=
				\tensor{S}{^d_e}\tensor{T}{^{ab}_c}
				=
				\tensor{(S\otimes T)}{^{abc}_{de}}
			\]
			Since a tensor product makes a new tensor, it's also very easy to show.
			\[
				\tensor{T}{^{ab}_{c}}
				\tensor{S}{^d_e} 
				=
				\tensor{S}{^d_e}\tensor{T}{^{ab}_c}
				=
				\tensor{U}{^{abc}_{de}}
			\]
			From this, it is immediately clear that $U$ is a $(3,2)$ tensor. 

			It does not show the components, but it presents tensors in a more convenient way, similar 
			to the index notation that old physicists used.

			They then \textit{defined} that if a tensor symbol uses \textit{Greek} indices, then
			it is the components of a tensor, not the tensor itself. For example
			\[
				\tensor{T}{^\mu_{\nu\sigma}}
			\]
			are the tensor components. You can then combine it with the basis $p_\mu$, $dx^\nu$, and 
			others, so that it qualifies as a bona fide tensor.
			\[
				\tensor{T}{^\mu_{\nu\sigma}}\p_\mu \otimes dx^\nu \otimes dx^\sigma
			\]
			Therefore, a tensor can be represented as 
			\[
				\tensor{T}{^{a}_{bc}}
				=
				\tensor{T}{^\mu_{\nu\sigma}}\p_\mu\otimes dx^\nu\otimes dx^\sigma
			\]
			And a tensor acting on covectors and vectors can be easily represented as 
			\[
				T[\omega,\xi]=
				\tensor{T}{^{ab}_c}\omega_a \xi_b
			\]
			It is then immediately clear that $\omega$ and $\xi$ are covectors, and $T$ is
			lacking a vector input, since it is a $(2,1)$ tensor. You do not immediately
			get that information in index-free notation, most have to be explicitly stated, and you
			have to rely on context to understand things. It is also clear that it is a tensor \textit{acting} on other tensors,
			and not just components, and the indices indicate the slots being filled. There is no ambiguity, but it's easy to use. You therefore get both the convenience
			of index notation, and the rigour (should you desire) of the index-free notation.
			But tensor manipulation and operations itself are easier than the latter, but clearer than the former.
			I\footnote{Sir Vega} myself trained 
			as a relativist and took my PhD and Post Doc under those schools, my adviser studied
			under Penrose, so I also adapt this notation.

			One major drawback of this is again, if you work with higher dimensional tensors,
			you will run out of letters very quickly.
	\end{enumerate}

	Moving forward, we will use this abstract index notation to do our calculations.
	You may ask, ``Why teach index free notation at all''? It's that you can understand
	all the rigour of the mathematics before you can be free to break it and make heuristics 
	of your own. Index free notation is cumbersome, but it is exact, which is useful.
	Furthermore, if I did not teach you index-free notation, you would have a hard time understanding
	mathematical papers and books teaching differential geometry from a pure math perspective. 
	It would take some time and difficulty before you fully grasp what they mean, and there are very good, 
	comprehensive pure math differential geometry text books that are written purely under index-free
	notation. 
	You would not gain anything from those resources if you do not understand the notation they use.
	We learned it first, so that we know the foundations of what really it is we are doing.
	When we have truly leard all the rules, then we can start being lazy and do things
	in the abstract way we desire. However, never forget that the root of the notation you are using,
	which is the index-free rigourous definition. If you take your notation and run with it, 
	without appreciating the true meaning underneath it all, you will have a hard time extending
	your mathematics beyond where those assumptions your notation makes.
	

\subsubsection{Metric Tensor}
\subsubsection{Inner Product}
